\documentclass[12pt,a4paper]{article}

% ── Codificación y lenguaje ──────────────────────────────────────────────────
\usepackage[utf8]{inputenc}
\usepackage[T1]{fontenc}
\usepackage[spanish]{babel}

% ── Márgenes ────────────────────────────────────────────────────────────────
\usepackage[top=2.5cm, bottom=2.5cm, left=3cm, right=2.5cm]{geometry}

% ── Tipografía y formato ─────────────────────────────────────────────────────
\usepackage{times}
\usepackage{setspace}
\usepackage{parskip}
\setlength{\parindent}{0pt}

% ── Gráficos y colores ───────────────────────────────────────────────────────
\usepackage{graphicx}
\usepackage{xcolor}
\usepackage{float}

% ── Código fuente ────────────────────────────────────────────────────────────
\usepackage{listings}
\lstset{
  language=C++,
  basicstyle=\ttfamily\small,
  keywordstyle=\color{blue}\bfseries,
  commentstyle=\color{gray}\itshape,
  stringstyle=\color{red!70!black},
  numbers=left,
  numberstyle=\tiny\color{gray},
  numbersep=8pt,
  frame=single,
  breaklines=true,
  tabsize=4,
  captionpos=b,
  showstringspaces=false,
  backgroundcolor=\color{gray!8},
}

% ── Matemáticas ──────────────────────────────────────────────────────────────
\usepackage{amsmath}
\usepackage{amssymb}

% ── Tablas ───────────────────────────────────────────────────────────────────
\usepackage{booktabs}
\usepackage{array}
\usepackage{multirow}

% ── Cabecera y pie de página ─────────────────────────────────────────────────
\usepackage{fancyhdr}
\pagestyle{fancy}
\fancyhf{}
\fancyhead[L]{\small Computación Gráfica I -- Práctica N° 2}
\fancyhead[R]{\small UNSAAC}
\fancyfoot[C]{\thepage}
\renewcommand{\headrulewidth}{0.4pt}

% ── Hipervínculos ────────────────────────────────────────────────────────────
\usepackage[hidelinks]{hyperref}

% ════════════════════════════════════════════════════════════════════════════
\begin{document}

% ── Portada ──────────────────────────────────────────────────────────────────
\begin{titlepage}
  \centering
  \vspace*{0.5cm}

  {\large \textbf{Universidad Nacional de San Antonio Abad del Cusco}}\\[0.3cm]
  {\normalsize Facultad de Ingeniería Eléctrica, Electrónica, Informática y Mecánica}\\[0.1cm]
  {\normalsize Departamento Académico de Ingeniería Informática y de Sistemas}\\[1.5cm]

  \rule{\linewidth}{0.5mm}\\[0.4cm]
  {\LARGE \textbf{COMPUTACIÓN GRÁFICA I}}\\[0.3cm]
  {\Large \textbf{Práctica N° 2}}\\[0.2cm]
  {\large \textbf{Algoritmos para Trazo de Líneas}}\\[0.2cm]
  \rule{\linewidth}{0.5mm}\\[1.8cm]

  \begin{minipage}{0.48\textwidth}
    \begin{flushleft}
      \textbf{Docente:}\\
      Dr. Hans Harley Ccacyahuillca Bejar
    \end{flushleft}
  \end{minipage}
  \hfill
  \begin{minipage}{0.48\textwidth}
    \begin{flushright}
      \textbf{Estudiante:}\\
      Efrain Vitorino Marin\\
      \textbf{Código:} 160337
    \end{flushright}
  \end{minipage}

  \vfill
  {\normalsize Cusco, \today}
\end{titlepage}

% ── Tabla de contenidos ───────────────────────────────────────────────────────
\tableofcontents
\newpage

% ════════════════════════════════════════════════════════════════════════════
\section{Objetivos}

\begin{itemize}
  \item Conocer los algoritmos clásicos de trazo de líneas en computación gráfica.
  \item Comprender la base matemática de cada algoritmo y sus diferencias en eficiencia.
  \item Implementar los algoritmos de trazo de líneas utilizando OpenGL y GLUT en C++.
  \item Desarrollar el ejercicio propuesto: dibujar 5 líneas al azar y señalar sus puntos de intersección.
\end{itemize}

% ════════════════════════════════════════════════════════════════════════════
\section{Base Teórica}

\subsection{Ecuación de la Recta}

La ecuación cartesiana de una línea recta es:
\begin{equation}
  y = mx + b
\end{equation}

donde $m$ es la pendiente y $b$ la intersección con el eje $y$. Dados dos puntos
$(x_1, y_1)$ y $(x_2, y_2)$, se calculan como:

\begin{equation}
  m = \frac{y_2 - y_1}{x_2 - x_1}, \qquad b = y_1 - m\,x_1
\end{equation}

Los intervalos incrementales son:
\begin{equation}
  \Delta y = m\,\Delta x, \qquad \Delta x = \frac{\Delta y}{m}
\end{equation}

Si $|m| < 1$ la línea es más horizontal; si $|m| > 1$ es más vertical.

% ─────────────────────────────────────────────────────────────────────────────
\subsection{Algoritmo Pendiente-Intersección}

Es el método más directo: para cada valor entero de $x$ se evalúa la ecuación
$y = mx + b$ y se redondea al entero más cercano.

\textbf{Ventaja:} sencillo de implementar.\\
\textbf{Desventaja:} requiere una multiplicación y un redondeo flotante en cada paso.

% ─────────────────────────────────────────────────────────────────────────────
\subsection{Algoritmo DDA (Digital Differential Analyzer)}

En lugar de evaluar $y = mx + b$ en cada paso, acumula el incremento:
\begin{equation}
  y_{k+1} = y_k + m
\end{equation}

Se evita la multiplicación del bucle, pero aún utiliza aritmética de punto flotante
y un redondeo por iteración.

% ─────────────────────────────────────────────────────────────────────────────
\subsection{Algoritmo de Bresenham}

Trabaja exclusivamente con aritmética entera. Mantiene una variable de error $e$
que representa dos veces la distancia entre el píxel elegido y la línea ideal.
En cada paso:

\begin{equation}
  e \leftarrow e + \Delta y
\end{equation}

Si $2e \geq \Delta x$, se incrementa $y$ y se corrige el error:
$e \leftarrow e - \Delta x$.

El algoritmo se generaliza a los 8 octantes mediante:
\begin{itemize}
  \item \textbf{Transposición} (\textit{steep}): intercambia $x$ e $y$ cuando $|\Delta y| > |\Delta x|$.
  \item \textbf{Paso negativo} (\texttt{dy\_step}): permite pendientes negativas.
\end{itemize}

% ─────────────────────────────────────────────────────────────────────────────
\subsection{Algoritmo de Punto Medio}

Variante incremental de enteros para el primer octante ($0 < m \leq 1$).
Usa una variable de decisión $d_m$ inicializada como:
\begin{equation}
  d_m = 2\,\Delta y - \Delta x
\end{equation}

En cada paso, si $d_m \leq 0$ el punto medio está por debajo de la línea y se
mantiene $y$; si $d_m > 0$ se incrementa $y$:
\begin{equation}
  d_m \leftarrow
  \begin{cases}
    d_m + 2\,\Delta y & \text{si } d_m \leq 0\\
    d_m + 2\,\Delta y - 2\,\Delta x & \text{si } d_m > 0
  \end{cases}
\end{equation}

% ════════════════════════════════════════════════════════════════════════════
\section{Implementación}

\subsection{Estructura del Proyecto}

\begin{verbatim}
laboComGraficaGuia2/
├── algoritmos_lineas.cpp    # Prueba de los 4 algoritmos
├── ejercicio_propuesto.cpp  # 5 líneas + detección de intersecciones
├── Makefile
└── documentacion/
    └── informe.tex
\end{verbatim}

\subsection{Función Auxiliar \texttt{plot}}

Todos los algoritmos utilizan la función \texttt{plot(x, y)} que dibuja un
píxel en las coordenadas enteras indicadas mediante OpenGL:

\begin{lstlisting}[caption={Función plot en OpenGL}]
void plot(int ix, int iy) {
    glBegin(GL_POINTS);
        glVertex2i(ix, iy);
    glEnd();
}
\end{lstlisting}

% ─────────────────────────────────────────────────────────────────────────────
\subsection{Algoritmo Pendiente-Intersección}

\begin{lstlisting}[caption={Implementación Pendiente-Intersección}]
void PendienteInterseccion(int x0, int y0, int x1, int y1) {
    if (x0 > x1) { std::swap(x0, x1); std::swap(y0, y1); }
    float m = (float)(y1 - y0) / (x1 - x0);
    float b = (float)y0 - m * (float)x0;
    for (int x = x0; x <= x1; x++) {
        plot(x, (int)roundf(m * (float)x + b));
    }
}
\end{lstlisting}

% ─────────────────────────────────────────────────────────────────────────────
\subsection{Algoritmo DDA}

\begin{lstlisting}[caption={Implementación DDA}]
void DDA(int x0, int y0, int x1, int y1) {
    if (x0 > x1) { std::swap(x0, x1); std::swap(y0, y1); }
    float m = (float)(y1 - y0) / (x1 - x0);
    float y = (float)y0;
    for (int x = x0; x <= x1; x++) {
        plot(x, (int)roundf(y));
        y += m;
    }
}
\end{lstlisting}

% ─────────────────────────────────────────────────────────────────────────────
\subsection{Algoritmo de Bresenham}

\begin{lstlisting}[caption={Implementación Bresenham (todos los octantes)}]
void Bresenham(int x0, int y0, int x1, int y1) {
    bool steep = std::abs(y1 - y0) > std::abs(x1 - x0);
    if (steep) { std::swap(x0, y0); std::swap(x1, y1); }
    if (x0 > x1) { std::swap(x0, x1); std::swap(y0, y1); }

    int dy_step = (y0 < y1) ? 1 : -1;
    int delta_x = x1 - x0;
    int delta_y = std::abs(y1 - y0);
    int y       = y0;
    int error   = 0;

    for (int x = x0; x <= x1; x++) {
        if (steep) plot(y, x); else plot(x, y);
        error += delta_y;
        if (2 * error >= delta_x) {
            y     += dy_step;
            error -= delta_x;
        }
    }
}
\end{lstlisting}

% ─────────────────────────────────────────────────────────────────────────────
\subsection{Algoritmo de Punto Medio}

\begin{lstlisting}[caption={Implementación Punto Medio (primer octante)}]
void PtoMedio(int x0, int y0, int x1, int y1) {
    int delta_x = x1 - x0;
    int delta_y = y1 - y0;
    int dm      = 2 * delta_y - delta_x;
    int dm_p1   = 2 * delta_y;
    int dm_p2   = 2 * delta_y - 2 * delta_x;
    int y       = y0;

    for (int x = x0; x <= x1; x++) {
        plot(x, y);
        if (dm <= 0) {
            dm += dm_p1;
        } else {
            dm += dm_p2;
            y++;
        }
    }
}
\end{lstlisting}

% ════════════════════════════════════════════════════════════════════════════
\section{Ejercicio Propuesto}

\subsection{Descripción}

El ejercicio consiste en:
\begin{enumerate}
  \item Dibujar \textbf{5 líneas al azar} usando el algoritmo de Bresenham.
  \item Calcular y \textbf{señalar todos los puntos de intersección} entre ellas.
\end{enumerate}

\subsection{Detección de Intersecciones}

Se utiliza la \textbf{forma paramétrica} de dos segmentos:
\begin{align}
  P(t) &= A + t\,(B - A), \quad t \in [0,1]\\
  Q(s) &= C + s\,(D - C), \quad s \in [0,1]
\end{align}

Igualando $P(t) = Q(s)$ se obtiene el sistema:
\begin{equation}
  D = \Delta x_2\,\Delta y_1 - \Delta x_1\,\Delta y_2
\end{equation}
\begin{equation}
  t = \frac{\Delta x_2\,(C_y - A_y) - \Delta y_2\,(C_x - A_x)}{D}, \qquad
  s = \frac{\Delta x_1\,(C_y - A_y) - \Delta y_1\,(C_x - A_x)}{D}
\end{equation}

Hay intersección de segmento cuando $D \neq 0$, $t \in [0,1]$ y $s \in [0,1]$.
Con 5 líneas existen $\binom{5}{2} = 10$ pares posibles a verificar.

\subsection{Implementación de la Detección}

\begin{lstlisting}[caption={Función de intersección paramétrica}]
bool segmentIntersect(const Line& a, const Line& b, Point2D& pt) {
    float dx1 = (float)(a.x1 - a.x0), dy1 = (float)(a.y1 - a.y0);
    float dx2 = (float)(b.x1 - b.x0), dy2 = (float)(b.y1 - b.y0);

    float D = dx2 * dy1 - dx1 * dy2;
    if (fabsf(D) < 1e-6f) return false;  // paralelos

    float cx = (float)(b.x0 - a.x0);
    float cy = (float)(b.y0 - a.y0);

    float t = (dx2 * cy - dy2 * cx) / D;
    float s = (dx1 * cy - dy1 * cx) / D;

    if (t >= 0.0f && t <= 1.0f && s >= 0.0f && s <= 1.0f) {
        pt.x = (float)a.x0 + t * dx1;
        pt.y = (float)a.y0 + t * dy1;
        return true;
    }
    return false;
}
\end{lstlisting}

\subsection{Líneas Utilizadas}

\begin{table}[H]
  \centering
  \begin{tabular}{ccccc}
    \toprule
    \textbf{Línea} & \textbf{x0} & \textbf{y0} & \textbf{x1} & \textbf{y1} \\
    \midrule
    0 &  20 & 380 & 490 & 250 \\
    1 & 100 & 490 & 400 &  10 \\
    2 &  10 & 140 & 490 & 360 \\
    3 & 240 & 490 & 310 &  10 \\
    4 & 480 & 450 &  10 & 170 \\
    \bottomrule
  \end{tabular}
  \caption{Coordenadas de las 5 líneas del ejercicio propuesto}
\end{table}

% ════════════════════════════════════════════════════════════════════════════
\section{Compilación y Ejecución}

\subsection{Dependencias}

\begin{itemize}
  \item Compilador \texttt{g++} con soporte C++11.
  \item Librería \texttt{freeglut3-dev} (OpenGL/GLUT).
\end{itemize}

Instalación en Ubuntu/Debian:
\begin{verbatim}
sudo apt install g++ make freeglut3-dev
\end{verbatim}

\subsection{Comandos}

\begin{verbatim}
# Compilar ambos programas
make

# Ejecutar prueba de los 4 algoritmos
./algoritmos_lineas

# Ejecutar ejercicio propuesto
./ejercicio_propuesto
\end{verbatim}

% ════════════════════════════════════════════════════════════════════════════
\section{Resultados}

\subsection{Algoritmos de Trazado}

El programa \texttt{algoritmos\_lineas} muestra una ventana de 500×500 px con
cuatro líneas de prueba, una por cada algoritmo, diferenciadas por color:

\begin{table}[H]
  \centering
  \begin{tabular}{llll}
    \toprule
    \textbf{Color}  & \textbf{Algoritmo}        & \textbf{Tipo aritmética} & \textbf{Octantes} \\
    \midrule
    Rojo            & Pendiente-Intersección    & Flotante                 & 1° -- 4°          \\
    Azul            & DDA                       & Flotante (incremental)   & 1° -- 4°          \\
    Verde           & Bresenham                 & Entera                   & Todos             \\
    Magenta         & Punto Medio               & Entera (incremental)     & 1°                \\
    \bottomrule
  \end{tabular}
  \caption{Comparación de los algoritmos implementados}
\end{table}

\subsection{Ejercicio Propuesto}

El programa \texttt{ejercicio\_propuesto} muestra las 5 líneas en azul y marca
los puntos de intersección con cuadrados rojos de $9\times9$ px.

% ════════════════════════════════════════════════════════════════════════════
\section{Conclusiones}

\begin{enumerate}
  \item El algoritmo de \textbf{Pendiente-Intersección} es el más intuitivo pero
        el menos eficiente al requerir una multiplicación flotante por píxel.

  \item El \textbf{DDA} mejora la eficiencia al reemplazar la multiplicación por
        una suma acumulativa, aunque sigue dependiendo de aritmética flotante.

  \item \textbf{Bresenham} es el algoritmo más eficiente al usar solo sumas y
        restas enteras. Además, su versión generalizada cubre todos los octantes,
        lo que lo hace el más robusto de los tres.

  \item El \textbf{Punto Medio} ofrece la misma eficiencia entera que Bresenham
        pero en su forma básica se restringe al primer octante, requiriendo
        transformaciones adicionales para generalizarse.

  \item La detección de intersecciones mediante la forma paramétrica es precisa
        y eficiente, evaluando exactamente los 10 pares posibles entre 5 segmentos.
\end{enumerate}

% ════════════════════════════════════════════════════════════════════════════
\section*{Bibliografía}
\addcontentsline{toc}{section}{Bibliografía}

\begin{itemize}
  \item Donald Hearn, M. Pauline Baker; \textit{Gráficas por Computadora con OpenGL};
        Pearson Prentice Hall; 3° edición; Madrid 2006.

  \item Rimon Elias; \textit{Digital Media: A Problem-Solving Approach for Computer
        Graphics}; eBook. New York.

  \item Sumanta Guha; \textit{Computer Graphics Through OpenGL}; Third Edition;
        Taylor \& Francis Group. 2019.
\end{itemize}

\end{document}
